% math-model.tex
%
% Mathematical documentation of the four-wire rectangular IVR-EN optimal
% power flow implemented in BMOPFTools.jl (ext/BMOPFOpfExt).
%
% Intended for inclusion in the IEEE PES Task Force specification.
% Compile standalone with pdflatex / lualatex.

\documentclass[11pt,a4paper]{article}

\usepackage{amsmath,amssymb,amsthm}
\usepackage{mathtools}          % dcases, etc.
\usepackage{bm}                 % bold math
\usepackage{booktabs}
\usepackage{array}
\usepackage{xcolor}
\usepackage[hidelinks]{hyperref}
\usepackage{geometry}
\geometry{margin=2.5cm}

\newcommand{\R}{\mathbb{R}}
\newcommand{\C}{\mathbb{C}}
\newcommand{\mc}[1]{\mathcal{#1}}
\newcommand{\vr}[2]{v^r_{#1,#2}}
\newcommand{\vi}[2]{v^i_{#1,#2}}
\newcommand{\dvr}{\Delta v^r}
\newcommand{\dvi}{\Delta v^i}
\newcommand{\Zmat}{\mathbf{Z}}
\newcommand{\Rmat}{\mathbf{R}}
\newcommand{\Xmat}{\mathbf{X}}

\title{\textbf{Four-Wire Rectangular Current--Voltage OPF}\\[4pt]
       \large Mathematical Model Reference --- BMOPFTools.jl}
\author{IEEE PES Task Force on Benchmark Models for\\
        Optimal Power Flow of Distribution Systems}
\date{\today}

% ============================================================
\begin{document}
\maketitle
\tableofcontents
\bigskip

% ============================================================
\section{Overview}

BMOPFTools implements a \emph{four-wire rectangular current--voltage} (IVR-EN)
optimal power flow (OPF) for unbalanced distribution networks.  The formulation
is a direct extension of the \texttt{IVRENPowerModel} from
PowerModelsDistribution~\cite{PowerModelsDistribution2022} to native BMOPF JSON
network data.

Key design choices:
\begin{itemize}
  \item \textbf{Rectangular coordinates throughout.}  Voltage and current
        variables are real/imaginary pairs $(v^r, v^i)$ and $(c^r, c^i)$ in
        SI units (V and A).  No polar transformation is applied.

  \item \textbf{Explicit neutral conductor.}  The neutral terminal is a
        first-class variable; its voltage floats and is determined jointly by
        the KCL equations and any grounding conditions declared on the bus.
        This is the ``EN'' (explicit neutral) aspect of the formulation.

  \item \textbf{Constant-power loads and generators.}  Active and reactive
        power are specified as fixed $P^\text{nom}$, $Q^\text{nom}$ (loads) or
        bounded intervals $[P^\text{min}, P^\text{max}]$,
        $[Q^\text{min}, Q^\text{max}]$ (generators).  The resulting power
        equations are \emph{bilinear} in the voltage and current variables,
        making the problem a non-convex NLP.

  \item \textbf{Shunt admittances are modelled.}  The $\pi$-model half-sections
        ($G_\text{fr/to}$, $B_\text{fr/to}$ on linecodes, scaled by line length)
        and standalone shunt objects ($G$/$B$ matrices, bus-connected to ground)
        contribute linear current terms to KCL.  Thermal current limits are
        enforced on the total current (series + shunt) at both line ends.

  \item \textbf{NLP solver.}  The model is built with JuMP and solved with
        Ipopt.  A feasibility-relaxed variant (\S\ref{sec:feasibility}) adds
        elastic slack currents so that the KCL system is always satisfiable,
        enabling diagnostics on infeasible data.
\end{itemize}

% ============================================================
\section{Sets and Indices}

\begin{center}
\renewcommand{\arraystretch}{1.25}
\begin{tabular}{@{}lll@{}}
\toprule
Symbol & Description & Defined on \\
\midrule
$\mc{B}$           & Buses & network \\
$\mc{T}_b$         & Terminals of bus $b$ & bus \\
$\mc{T}_b^\phi$    & Phase (non-neutral) terminals of bus $b$ & bus \\
$n_b$              & Neutral terminal name of bus $b$ (if any) & bus \\
$\mc{G}_\text{nd}$ & Grounded terminals: $(b,t)$ s.t.\ $v=0$ & bus \\
$\mc{L}$           & Lines & network \\
$\mc{S}$           & Switches & network \\
$\mc{D}$           & Loads & network \\
$\mc{G}$           & Generators & network \\
$\mc{N}$           & IBRs & network \\
$\mc{V}$           & Voltage sources (slack buses) & network \\
$\mc{X}^\text{YY}$ & Single-phase wye--wye transformers (\texttt{single\_phase}) & network \\
$\mc{X}^\text{CT}$ & Center-tap split-phase transformers (\texttt{center\_tap}) & network \\
$\mc{X}^\text{Yd}$ & Wye--delta transformers & network \\
$\mc{X}^\text{Dy}$ & Delta--wye transformers & network \\
\bottomrule
\end{tabular}
\end{center}

The neutral terminal $n_b$ is identified in priority order:
(1)~the explicit \texttt{neutral\_terminal} field on the bus,
(2)~a terminal whose name is \texttt{"n"} or \texttt{"N"} (case-insensitive).
If neither condition holds, all terminals of bus~$b$ are treated as phase
conductors.

% ============================================================
\section{Parameters}

\begin{center}
\renewcommand{\arraystretch}{1.25}
\begin{tabular}{@{}llll@{}}
\toprule
Symbol & Unit & Description & Source field \\
\midrule
$\Rmat_\ell, \Xmat_\ell$ & $\Omega$ & Total series resistance / reactance matrix of line $\ell$ & \texttt{linecode} $\times$ \texttt{length} \\
$I^\text{max}_{\ell,k}$ & A & Current magnitude limit, conductor $k$ of line $\ell$ & \texttt{linecode.i\_max} \\
$P^{d,\text{nom}}_{d,k}$, $Q^{d,\text{nom}}_{d,k}$ & W, var & Nominal active/reactive demand, phase $k$ & \texttt{load.p\_nom}, \texttt{q\_nom} \\
$P^{g,\text{min/max}}_{g,k}$, $Q^{g,\text{min/max}}_{g,k}$ & W, var & Generator active/reactive power limits & \texttt{generator.p\_min/max}, \texttt{q\_min/max} \\
$c^g_1$ & \$/W & Linear generation cost coefficient & \texttt{generator.cost} \\
$P^{n,\text{min/max}}_{n,k}$, $Q^{n,\text{min/max}}_{n,k}$ & W, var & IBR active/reactive power limits & \texttt{ibr.p\_min/max}, \texttt{q\_min/max} \\
$S^{\max}_{n,k}$ & VA & IBR apparent-power rating, phase $k$ & \texttt{ibr.s\_max} \\
$\mathrm{pf}_n$ & -- & Signed constant power factor (if set) & \texttt{control\_profile.power\_factor.pf} \\
$V^s_{v,k}$, $\theta^s_{v,k}$ & V, rad & Voltage source magnitude and angle at terminal $k$ & \texttt{voltage\_source.v\_magnitude/angle} \\
$v^b_\text{min}$, $v^b_\text{max}$ & V & Phase voltage magnitude limits at bus $b$ & \texttt{bus.v\_min}, \texttt{v\_max} \\
$N_x$ & -- & Transformer turns ratio $= V^\text{ref}_\text{fr} / V^\text{ref}_\text{to}$ & \texttt{transformer.v\_ref\_from/to} \\
$R_1, X_1$ & $\Omega$ & From-side winding resistance/reactance (HV base) & \texttt{transformer.r\_series\_from/x\_series\_from} \\
$R_2, X_2$ & $\Omega$ & To-side winding resistance/reactance (LV base) & \texttt{transformer.r\_series\_to/x\_series\_to} \\
$G_0$ & S & No-load (core-loss) conductance, placed at HV terminals & \texttt{transformer.g\_no\_load} \\
$B_0$ & S & No-load (magnetising) susceptance, placed at HV terminals & \texttt{transformer.b\_no\_load} \\
\bottomrule
\end{tabular}
\end{center}

The series impedance matrices $\Rmat_\ell$, $\Xmat_\ell \in \R^{n_c \times n_c}$
are constructed from the line's associated \texttt{linecode} (which stores
impedance per metre) multiplied by \texttt{length} in metres.  The matrices are
symmetric by the Carson/primitive impedance construction.

% ============================================================
\section{Decision Variables}

All variables are real-valued and in SI units.

\begin{center}
\renewcommand{\arraystretch}{1.25}
\begin{tabular}{@{}llll@{}}
\toprule
Variable & Index & Unit & Description \\
\midrule
$v^r_{b,t}$, $v^i_{b,t}$   & $(b,t)\in\mc{B}\times\mc{T}_b$ & V & Rectangular voltage \\
$c^r_{\ell,k}$, $c^i_{\ell,k}$ & $\ell\in\mc{L},\,k=1\ldots n_\ell$ & A & Series current (from-side), conductor $k$ \\
$\tilde{c}^r_{\ell,k}$, $\tilde{c}^i_{\ell,k}$ & same & A & Series current (to-side), conductor $k$ \\
$c^{r,s}_{s,k}$, $c^{i,s}_{s,k}$ & $s\in\mc{S},\,k=1\ldots n_s$ & A & Switch current, conductor $k$ \\
$c^{r,d}_{d,k}$, $c^{i,d}_{d,k}$ & $d\in\mc{D},\,k=1\ldots|\mc{T}_d^\phi|$ & A & Load current, phase conductor $k$ \\
$c^{r,g}_{g,k}$, $c^{i,g}_{g,k}$ & $g\in\mc{G},\,k=1\ldots|\mc{T}_g^\phi|$ & A & Generator current, phase conductor $k$ \\
$c^{r,n}_{n,k}$, $c^{i,n}_{n,k}$ & $n\in\mc{N},\,k=1\ldots n^\phi_n$ & A & IBR current, phase conductor $k$ \\
$c^{r,v}_{v,k}$, $c^{i,v}_{v,k}$ & $v\in\mc{V},\,k=1\ldots n_v$ & A & Voltage source slack current \\
$c^{r,x}_{x,\sigma,k}$, $c^{i,x}_{x,\sigma,k}$ & $x\in\mc{X},\,\sigma\in\{\text{fr},\text{to}\}$ & A & Transformer winding current \\
\bottomrule
\end{tabular}
\end{center}

Generator and load current variables are indexed over the \emph{phase}
conductors only; the neutral return current is not an independent variable but
is determined implicitly by KCL.

% ============================================================
\section{Objective Function}

The objective minimises total active-power generation cost:
%
\begin{equation}
  \min \sum_{g\in\mc{G}} \sum_{k=1}^{|\mc{T}_g^\phi|}
    c^g_1 \cdot P_{g,k},
\end{equation}
%
where the per-phase generated power $P_{g,k}$ is a bilinear expression in
the decision variables (see \S\ref{sec:generators}).

The cost field on a generator may alternatively be specified as a three-vector
$[c_2, c_1, c_0]$.  The constant $c_0$ does not affect the optimum and is
omitted.  The quadratic coefficient $c_2$ is currently not supported (a warning
is issued); only the linear term $c_1 P$ is included.

% ============================================================
\section{Constraints}

\subsection{Grounding}

Terminals declared as perfectly grounded are fixed to zero potential:
%
\begin{equation}
  \vr{b}{t} = 0,\quad \vi{b}{t} = 0, \qquad \forall\,(b,t)\in\mc{G}_\text{nd}.
\end{equation}

\subsection{Voltage Sources (Slack Bus)}

A voltage source $v\in\mc{V}$ fixes the rectangular voltage at each of its
terminals and injects an unconstrained slack current into the KCL system:
%
\begin{align}
  \vr{b}{t_k} &= V^s_{v,k}\cos\theta^s_{v,k}, \\
  \vi{b}{t_k} &= V^s_{v,k}\sin\theta^s_{v,k},
\end{align}
%
for each terminal $t_k$ in the source's \texttt{terminal\_map}.  The slack
current $(c^{r,v}_{v,k}, c^{i,v}_{v,k})$ is a free variable and absorbs
whatever power balance the rest of the network requires, making the source bus
the power-flow reference bus.

\subsection{Voltage Magnitude Bounds}

Operational voltage limits are enforced at every ungrounded, non-source,
non-neutral phase terminal:
%
\begin{equation}
  \left(v^b_\text{min}\right)^2
  \;\leq\;
  \left(\vr{b}{t}\right)^2 + \left(\vi{b}{t}\right)^2
  \;\leq\;
  \left(v^b_\text{max}\right)^2,
  \qquad \forall\,b\in\mc{B},\; t\in\mc{T}_b^\phi.
\end{equation}
%
The neutral terminal voltage is excluded from voltage bounds; it is determined
by physics (KCL) and may float within whatever range the network equations
permit.  The squared form avoids introducing a non-smooth square root.

\subsection{Lines}
\label{sec:lines}

\subsubsection*{KVL (series voltage drop)}

For line $\ell$ with $n_\ell$ conductors, total series impedance
$\Zmat_\ell = \Rmat_\ell + j\Xmat_\ell \in \C^{n_\ell\times n_\ell}$, and
terminal-map arrays $\mathbf{t}^\text{fr}_\ell$, $\mathbf{t}^\text{to}_\ell$:
%
\begin{align}
  \vr{b^\text{fr}}{t^\text{fr}_k}
  - \vr{b^\text{to}}{t^\text{to}_k}
  &= \sum_{j=1}^{n_\ell}
     \bigl(R_{\ell,kj}\,c^r_{\ell,j} - X_{\ell,kj}\,c^i_{\ell,j}\bigr),
  \label{eq:kvl-r} \\[4pt]
  \vi{b^\text{fr}}{t^\text{fr}_k}
  - \vi{b^\text{to}}{t^\text{to}_k}
  &= \sum_{j=1}^{n_\ell}
     \bigl(R_{\ell,kj}\,c^i_{\ell,j} + X_{\ell,kj}\,c^r_{\ell,j}\bigr),
  \label{eq:kvl-i}
\end{align}
%
for $k = 1,\ldots,n_\ell$.  The impedance matrix $\Zmat_\ell$ captures full
mutual coupling between conductors; the off-diagonal terms couple the voltage
drop on one conductor to the current in another.

\subsubsection*{Series current balance}

The to-side series current variable equals the negative of the from-side:
%
\begin{equation}
  \tilde{c}^r_{\ell,k} = -c^r_{\ell,k},\qquad
  \tilde{c}^i_{\ell,k} = -c^i_{\ell,k},\qquad
  k = 1,\ldots,n_\ell.
\end{equation}

\subsubsection*{$\pi$-model shunt currents}

The linecode $G_\text{fr/to}$, $B_\text{fr/to}$ fields (S/m, scaled by line
length) define shunt admittance half-sections at each end.  The resulting shunt
current at bus~$b$, conductor~$k$ is linear in the voltage variables:
%
\begin{equation}
  I^{\text{sh},r}_{k}(b)
  = \sum_j \bigl(G_{kj}\,v^r_{b,t_j} - B_{kj}\,v^i_{b,t_j}\bigr),\qquad
  I^{\text{sh},i}_{k}(b)
  = \sum_j \bigl(G_{kj}\,v^i_{b,t_j} + B_{kj}\,v^r_{b,t_j}\bigr).
  \label{eq:ish}
\end{equation}
%
No new JuMP variables are introduced; the expressions are accumulated directly
into the KCL terms.

\subsubsection*{Current magnitude limits}

Thermal limits are enforced on the \emph{total} current (series + shunt) at
both ends of the line.  When the linecode specifies $I^\text{max}_{\ell,k}$:
%
\begin{align}
  \bigl(c^r_{\ell,k} + I^{\text{sh},r}_k(b^\text{fr})\bigr)^2
  + \bigl(c^i_{\ell,k} + I^{\text{sh},i}_k(b^\text{fr})\bigr)^2
  &\;\leq\; \bigl(I^\text{max}_{\ell,k}\bigr)^2,
  \label{eq:i-lim-fr} \\[4pt]
  \bigl(\tilde{c}^r_{\ell,k} + I^{\text{sh},r}_k(b^\text{to})\bigr)^2
  + \bigl(\tilde{c}^i_{\ell,k} + I^{\text{sh},i}_k(b^\text{to})\bigr)^2
  &\;\leq\; \bigl(I^\text{max}_{\ell,k}\bigr)^2,
  \label{eq:i-lim-to}
\end{align}
%
for $k = 1,\ldots,n_\ell$.  When shunt fields are absent or zero, these
collapse to the familiar series-only constraint at both ends.

\subsubsection*{KCL contributions}

At bus~$b^\text{fr}$, terminal $t^\text{fr}_k$: the series current
$c^r_{\ell,k}$ and the from-end shunt current $I^{\text{sh},r}_k(b^\text{fr})$
both leave the bus, contributing $-(c^r_{\ell,k} + I^{\text{sh},r}_k(b^\text{fr}))$.
At bus~$b^\text{to}$, terminal $t^\text{to}_k$: the to-side series current
$\tilde{c}^r_{\ell,k} = -c^r_{\ell,k}$ and the to-end shunt current both leave
the bus, contributing $-(\tilde{c}^r_{\ell,k} + I^{\text{sh},r}_k(b^\text{to}))
= c^r_{\ell,k} - I^{\text{sh},r}_k(b^\text{to})$.

\subsection{Switches}

A \emph{closed} switch enforces zero voltage drop across its terminals:
%
\begin{equation}
  \vr{b^\text{fr}}{t^\text{fr}_k} = \vr{b^\text{to}}{t^\text{to}_k},\qquad
  \vi{b^\text{fr}}{t^\text{fr}_k} = \vi{b^\text{to}}{t^\text{to}_k},\qquad
  k = 1,\ldots,n_s.
\end{equation}
%
An \emph{open} switch has its current variables fixed to zero
$(c^{r,s}_{s,k} = c^{i,s}_{s,k} = 0)$; no voltage coupling is imposed, so
the two buses are electrically disconnected at that conductor.

Switch currents contribute to KCL with sign convention: current $c^{r,s}_{s,k}$
leaves $b^\text{fr}$ and enters $b^\text{to}$.

\subsection{Standalone Shunts}
\label{sec:shunts}

A \texttt{shunt} object represents an admittance matrix
$\mathbf{Y}^{sh} = \mathbf{G}^{sh} + j\mathbf{B}^{sh}$ (S) connected between
a set of bus terminals and ground.  The current drawn from terminal~$k$ is
linear in the voltage variables (no new JuMP variables):
%
\begin{equation}
  I^{\text{sh},r}_{k}
  = \sum_j \bigl(G^{sh}_{kj}\,v^r_{b,t_j} - B^{sh}_{kj}\,v^i_{b,t_j}\bigr),\qquad
  I^{\text{sh},i}_{k}
  = \sum_j \bigl(G^{sh}_{kj}\,v^i_{b,t_j} + B^{sh}_{kj}\,v^r_{b,t_j}\bigr).
\end{equation}
%
KCL contribution: $-I^{\text{sh},r}_k$ at terminal~$t_k$ (current leaves the
bus to ground).  Grounded terminals are absent from the voltage variable dict
and contribute nothing.

\subsection{Loads}
\label{sec:loads}

The load model is \emph{constant power}: active and reactive power at each
phase conductor are fixed parameters $P^{d,\text{nom}}_{d,k}$ and
$Q^{d,\text{nom}}_{d,k}$.  The relationship between power, voltage, and current
is bilinear in rectangular coordinates.

\subsubsection*{Wye / Single-phase loads}

Let $t^\phi_k$ denote the $k$-th phase terminal of load $d$ and $n_b$ the
neutral terminal of bus~$b$ (or implicitly zero if absent).  Define the
phase-to-neutral voltage drop:
%
\begin{equation}
  \dvr_k = \vr{b}{t^\phi_k} - \vr{b}{n_b},
  \qquad
  \dvi_k = \vi{b}{t^\phi_k} - \vi{b}{n_b}.
\end{equation}
%
The constant-power equations for phase $k$ are:
%
\begin{align}
  \dvr_k\,c^{r,d}_{d,k} + \dvi_k\,c^{i,d}_{d,k} &= P^{d,\text{nom}}_{d,k},
  \label{eq:load-p}\\
  \dvi_k\,c^{r,d}_{d,k} - \dvr_k\,c^{i,d}_{d,k} &= Q^{d,\text{nom}}_{d,k}.
  \label{eq:load-q}
\end{align}
%
These are the real and imaginary parts of $S = V \cdot I^*$ in rectangular form.
The neutral return current is not an explicit variable; it appears as the
negative sum of phase currents in KCL at terminal $n_b$.

\subsubsection*{Delta loads}

For a delta-connected load with $n_c$ terminals, element $k$ is connected
positive-to-negative between terminal $t_k$ and $t_{k^+}$ where
$k^+ = (k \bmod n_c) + 1$ (cyclic):
%
\begin{equation}
  \dvr_k = \vr{b}{t_k} - \vr{b}{t_{k^+}},
  \qquad
  \dvi_k = \vi{b}{t_k} - \vi{b}{t_{k^+}}.
\end{equation}
%
The power equations (\ref{eq:load-p})--(\ref{eq:load-q}) apply with these
line-to-line voltage drops.  KCL contributions are: $-c^{r,d}_{d,k}$ at
terminal $t_k$ and $+c^{r,d}_{d,k}$ at terminal $t_{k^+}$.

\subsection{Generators}
\label{sec:generators}

Generators follow the same bilinear power expressions as wye-connected loads,
but with power \emph{injected} rather than consumed, and with bounds instead of
equalities:
%
\begin{align}
  P^{g,\text{min}}_{g,k}
  \;\leq\; \dvr_k\,c^{r,g}_{g,k} + \dvi_k\,c^{i,g}_{g,k}
  &\;\leq\; P^{g,\text{max}}_{g,k},
  \label{eq:gen-p}\\
  Q^{g,\text{min}}_{g,k}
  \;\leq\; \dvi_k\,c^{r,g}_{g,k} - \dvr_k\,c^{i,g}_{g,k}
  &\;\leq\; Q^{g,\text{max}}_{g,k}.
  \label{eq:gen-q}
\end{align}
%
The voltage drop $(\dvr_k, \dvi_k)$ is phase-to-neutral for WYE generators and
line-to-line for DELTA generators, following the same convention as loads.
Generator currents contribute positively to KCL at the phase terminal and
negatively at the neutral terminal (current flows \emph{out of} the neutral).

\subsection{IBRs}
\label{sec:ibrs}

IBRs reuse the generator power model with the per-phase active and reactive
powers
%
\begin{align}
  P_{n,k} &= \dvr_k\,c^{r,n}_{n,k} + \dvi_k\,c^{i,n}_{n,k}, &
  Q_{n,k} &= \dvi_k\,c^{r,n}_{n,k} - \dvr_k\,c^{i,n}_{n,k}.
  \label{eq:inv-pq}
\end{align}
%
The voltage drop $(\dvr_k, \dvi_k)$ and the number of current variables
$n^\phi_n$ follow the inverter \texttt{topology}:
%
\begin{itemize}
  \item \textbf{\texttt{FOUR\_LEG}} --- phase-to-neutral,
        $\Delta v_k = v_{b,t_k} - v_{b,t_n}$, one current per phase conductor.
  \item \textbf{\texttt{THREE\_LEG}} --- line-to-line (delta),
        $\Delta v_k = v_{b,t_k} - v_{b,t_{k^+}}$ with cyclic index
        $k^+ = (k \bmod n) + 1$; no neutral current.
  \item \textbf{\texttt{SINGLE\_PHASE}} --- phase-to-reference,
        $\Delta v = v_{b,t_1} - v_{b,t_2}$, a single current.
\end{itemize}
%
Each phase is bounded in active power and, when an apparent-power rating is
given, by a quadratic apparent-power circle:
%
\begin{align}
  P^{n,\text{min}}_{n,k} \;\leq\; P_{n,k} &\;\leq\; P^{n,\text{max}}_{n,k},
  \label{eq:inv-p}\\
  P_{n,k}^2 + Q_{n,k}^2 &\;\leq\; \bigl(S^{\max}_{n,k}\bigr)^2.
  \label{eq:inv-s}
\end{align}
%
Reactive power is handled in one of two mutually exclusive ways.  By default
box bounds apply,
$Q^{n,\text{min}}_{n,k} \leq Q_{n,k} \leq Q^{n,\text{max}}_{n,k}$.  When the
IBR references a \texttt{control\_profile} with a signed power factor
$\mathrm{pf}_n$, $Q$ is instead coupled to $P$ by the exact equality
%
\begin{equation}
  \operatorname{sign}(\mathrm{pf}_n)\,Q_{n,k}
  + \tan\!\bigl(\arccos|\mathrm{pf}_n|\bigr)\,P_{n,k} = 0,
  \label{eq:inv-pf}
\end{equation}
%
with $\mathrm{pf}_n > 0$ lagging (absorbing VAr) and $\mathrm{pf}_n < 0$ leading
(injecting VAr).  IBR currents enter KCL with the generator sign convention
(injection positive at the phase terminal); for \texttt{FOUR\_LEG} the negated
phase current is added at the neutral terminal.

\subsection{Transformers}

All transformer constraints are linear in the voltage and current variables.
Apparent power and thermal limits on transformers are not currently modelled
as explicit constraints.

The turns ratio for all subtypes is
$N_x = V^\text{ref}_\text{fr} / V^\text{ref}_\text{to}$,
where both reference voltages are in SI volts.  For \texttt{center\_tap},
$V^\text{ref}_\text{to}$ is the \emph{per-leg} voltage (e.g.\ 120~V), not
the full secondary span (240~V).

\subsubsection*{Wye--wye per-phase (\texttt{single\_phase})}

The \texttt{single\_phase} subtype is modelled with a
$\Gamma$-equivalent circuit: series impedance referred to the from (HV) side,
and a shunt no-load admittance at the HV terminals.
Let $R_x = R_1 + N_x^2 R_2$ and $X_x = X_1 + N_x^2 X_2$ be the combined
series resistance and reactance referred to the HV side, where $R_1$, $X_1$
(\texttt{r\_series\_from}, \texttt{x\_series\_from}) are the HV winding
values in~$\Omega$ and $R_2$, $X_2$
(\texttt{r\_series\_to}, \texttt{x\_series\_to}) are the LV winding values
in~$\Omega$ on the LV base.

\paragraph{Voltage constraints.}
For each phase-pair index $k = 1,\ldots,n_c$:
%
\begin{align}
  \vr{b^\text{fr}}{t^\text{fr}_k} - N_x\,\vr{b^\text{to}}{t^\text{to}_k}
  &= R_x\,c^{r,x}_{x,\text{fr},k} - X_x\,c^{i,x}_{x,\text{fr},k},
  \label{eq:yy-v-r}\\
  \vi{b^\text{fr}}{t^\text{fr}_k} - N_x\,\vi{b^\text{to}}{t^\text{to}_k}
  &= R_x\,c^{i,x}_{x,\text{fr},k} + X_x\,c^{r,x}_{x,\text{fr},k}.
  \label{eq:yy-v-i}
\end{align}
%
When all series impedance fields are absent or zero the equations reduce to
the ideal constraint $V^\text{fr}_k = N_x V^\text{to}_k$.

\paragraph{Current coupling.}
The ideal-core coupling (power conservation) gives:
%
\begin{equation}
  N_x\,c^{r,x}_{x,\text{fr},k} + c^{r,x}_{x,\text{to},k} = 0,
  \qquad
  N_x\,c^{i,x}_{x,\text{fr},k} + c^{i,x}_{x,\text{to},k} = 0.
  \label{eq:yy-i}
\end{equation}

\paragraph{No-load shunt.}
The magnetising branch $Y_0 = G_0 + \mathrm{j}B_0$
(\texttt{g\_no\_load}, \texttt{b\_no\_load}, in~S) is placed at the HV
(from-side) terminals.  The total terminal current entering the HV bus
is the sum of the series current and the shunt:
%
\begin{equation}
  I^\text{fr,term}_{x,k} = c^{r,x}_{x,\text{fr},k}
    + G_0\,\vr{b^\text{fr}}{t^\text{fr}_k}
    - B_0\,\vi{b^\text{fr}}{t^\text{fr}_k}
    + \mathrm{j}\!\left(c^{i,x}_{x,\text{fr},k}
    + G_0\,\vi{b^\text{fr}}{t^\text{fr}_k}
    + B_0\,\vr{b^\text{fr}}{t^\text{fr}_k}\right).
  \label{eq:yy-shunt}
\end{equation}
%
This terminal current (not $c^{r,x}_{x,\text{fr},k}$ alone) enters the KCL
accumulator at terminal $t^\text{fr}_k$.

\subsubsection*{Center-tap split-phase (\texttt{center\_tap})}

The \texttt{center\_tap} subtype models the North American distribution
split-phase transformer: one HV winding drives two anti-series 120~V LV legs
sharing a center-tap neutral.

\paragraph{Terminal convention.}
\begin{itemize}
  \item HV (from) side: two terminals $[t^\text{ph}, t^\text{n}]$ —
        phase and neutral.
  \item LV (to) side: three terminals $[t_1, t_\text{n}, t_2]$ —
        leg-1, center-tap neutral, leg-2.
\end{itemize}
The center-tap terminal $t_\text{n}$ on the LV bus is typically
perfectly grounded (\texttt{perfectly\_grounded\_terminals}).

\paragraph{T-model physics.}
Unlike the two-winding $\Gamma$ model, the three-winding center-tap unit has
\emph{independent} LV winding branches.  Each leg carries a different current
under unbalanced loading, so lumping both secondary impedances into a single
referred-to-HV value would force identical leg voltages and fail to reproduce
load-imbalance effects.

The correct model is a T-circuit with a star node internal to the transformer:
\begin{itemize}
  \item HV branch: $Z_1 = R_1 + \mathrm{j}X_1$
        (\texttt{r\_series\_from}, \texttt{x\_series\_from})
        carries the total HV series current $I_s$.
  \item LV branch 1 (leg-1, $t_1 \to t_\text{n}$): $Z_2 = R_2 + \mathrm{j}X_2$
        (\texttt{r\_series\_to}, \texttt{x\_series\_to}) carries $I_{\ell_1}$.
  \item LV branch 2 (leg-2, $t_\text{n} \to t_2$): same $Z_2$, carries
        $I_{\ell_2}$.
\end{itemize}

Eliminating the internal star-node voltage gives the following constraint
for each leg $\ell \in \{1, 2\}$ with LV current $I_{\ell}$ and LV terminal
pair $(t_a, t_b)$:
%
\begin{align}
  \bigl(\vr{b^\text{fr}}{t^\text{ph}} - \vr{b^\text{fr}}{t^\text{n}}\bigr)
  - N_x\bigl(\vr{b^\text{to}}{t_a} - \vr{b^\text{to}}{t_b}\bigr)
  &= R_1\,c^{r,x}_{x,s} - X_1\,c^{i,x}_{x,s}
     - N_x\!\left(R_2\,c^{r,x}_{x,\ell} - X_2\,c^{i,x}_{x,\ell}\right),
  \label{eq:ct-v-r}\\
  \bigl(\vi{b^\text{fr}}{t^\text{ph}} - \vi{b^\text{fr}}{t^\text{n}}\bigr)
  - N_x\bigl(\vi{b^\text{to}}{t_a} - \vi{b^\text{to}}{t_b}\bigr)
  &= R_1\,c^{i,x}_{x,s} + X_1\,c^{r,x}_{x,s}
     - N_x\!\left(R_2\,c^{i,x}_{x,\ell} + X_2\,c^{r,x}_{x,\ell}\right),
  \label{eq:ct-v-i}
\end{align}
%
where $c^{r,x}_{x,s} = c^{r,x}_{x,\text{fr},1}$ is the HV series current
(variable index~1 on the \texttt{fr} side) and $c^{r,x}_{x,\ell}$ is the
corresponding LV leg current (variable index~1 or~3 on the \texttt{to} side).

The $-N_x Z_2 I_\ell$ sign arises from the T-model star-node elimination:
LV currents flow \emph{into} the transformer from the bus, hence opposite to
the direction through the winding branch from the star node.

\paragraph{Current coupling.}
Both LV currents flow into the transformer from their respective buses;
by power conservation at the ideal core:
%
\begin{equation}
  N_x\,c^{r,x}_{x,s} + c^{r,x}_{x,\ell_1} + c^{r,x}_{x,\ell_2} = 0,
  \qquad
  N_x\,c^{i,x}_{x,s} + c^{i,x}_{x,\ell_1} + c^{i,x}_{x,\ell_2} = 0.
  \label{eq:ct-i}
\end{equation}

\paragraph{Center-tap KCL.}
The center-tap terminal $t_\text{n}$ carries the imbalance current
$I_n = -I_{\ell_1} - I_{\ell_2}$:
%
\begin{equation}
  c^{r,x}_{x,n} + c^{r,x}_{x,\ell_1} + c^{r,x}_{x,\ell_2} = 0,
  \qquad
  c^{i,x}_{x,n} + c^{i,x}_{x,\ell_1} + c^{i,x}_{x,\ell_2} = 0,
  \label{eq:ct-n}
\end{equation}
%
where $c^{r,x}_{x,n} = c^{r,x}_{x,\text{to},2}$ (variable index~2 on the
\texttt{to} side).

\paragraph{No-load shunt.}
The shunt $G_0 + \mathrm{j}B_0$ is placed at the HV phase terminal
$t^\text{ph}$ (same convention as \texttt{single\_phase}, eq.~\eqref{eq:yy-shunt}).
The series current $I_s$ returns through the HV neutral $t^\text{n}$, so
the HV neutral KCL contribution is $+I_s$ (current entering the bus at $t^\text{n}$).
The shunt current is phase-to-ground and does not pass through $t^\text{n}$.

\paragraph{Impedance parameters from OpenDSS.}
For a 3-winding OpenDSS transformer with per-pair leakage values
\texttt{XHL}, \texttt{XLT}, \texttt{XHT} (all in~\%), the star-network
partitioning gives:
%
\begin{equation}
  X_1 = \frac{X_\text{HL} + X_\text{HT} - X_\text{LT}}{2}
        \cdot \frac{V_\text{hv}^2}{100\,S},
  \qquad
  X_2 = \frac{X_\text{HL} + X_\text{LT} - X_\text{HT}}{2}
        \cdot \frac{V_\text{lv}^2}{100\,S}.
  \label{eq:ct-star}
\end{equation}
%
For the symmetric case $X_\text{HT} = X_\text{HL}$ (both legs identical)
this simplifies to
$X_1 = (X_\text{HL} - X_\text{LT}/2)\cdot V_\text{hv}^2/(100S)$
and $X_2 = (X_\text{LT}/2)\cdot V_\text{lv}^2/(100S)$.
Note that $X_2 \neq X_\text{HL}/2$; using $X_\text{HL}/2$ for both windings
(the two-winding $\Gamma$ approximation) forces both leg voltages to be equal
under unbalanced loading and is incorrect for \texttt{center\_tap}.

The per-winding resistances map directly:
$R_1 = (\%r_1/100)\cdot V_\text{hv}^2/S$ and
$R_2 = (\%r_2/100)\cdot V_\text{lv}^2/S$.

\subsubsection*{Wye--delta (Yd) and delta--wye (Dy)}

The \texttt{wye\_delta} (Yd) and \texttt{delta\_wye} (Dy) subtypes share a
common formulation.  Let the subscripts ``wye'' and ``del'' refer to the wye-
and delta-connected windings respectively, with $n_\phi$ phase conductors on the
delta side.  The \emph{effective} per-winding turns ratio accounts for the
$\sqrt{3}$ factor introduced by the delta geometry:
%
\begin{equation}
  n_\text{eff} =
  \begin{dcases}
    \dfrac{\sqrt{3}}{N_x} & \text{Yd (wye is from-side)},\\[6pt]
    N_x\,\sqrt{3}         & \text{Dy (delta is from-side)}.
  \end{dcases}
\end{equation}
%
Here $N_x = V^\text{ref}_\text{fr} / V^\text{ref}_\text{to}$ and both reference
voltages are specified as phase-to-neutral equivalents.

\paragraph{Loss model (per-winding T).}
Following the OpenDSS reference model (reproduced by PowerModelsDistribution's
\texttt{eng2math} loss network), each winding carries its own series impedance
and a single no-load (core-loss) shunt $Y_0 = G_0 + \mathrm{j}B_0$
(\texttt{g\_no\_load}, \texttt{b\_no\_load}, in~S) sits at the from-side (HV)
phase terminals, phase-to-ground.  Write $R^\text{w}_x, X^\text{w}_x$
(\texttt{r\_series\_from}, \texttt{x\_series\_from}) for the wye/primary winding
branch and $R^\text{d}_x, X^\text{d}_x$
(\texttt{r\_series\_to}, \texttt{x\_series\_to}) for the delta/secondary branch,
both in~$\Omega$ on their own winding base.  A network supplying the legacy
single \texttt{r\_series}/\texttt{x\_series} is read as
$R^\text{w}_x = R_\text{series}$, $R^\text{d}_x = 0$ (likewise $X$),
recovering the ideal-delta behaviour.

\paragraph{Voltage constraints.}
For $k = 1,\ldots,n_\phi$ with $k^+ = (k \bmod n_\phi)+1$, the delta line-to-line
voltage equals the wye phase-to-neutral voltage scaled by $n_\text{eff}$, less
the series voltage drop across the winding branches:
%
\begin{align}
  \vr{\text{del}}{t_k} - \vr{\text{del}}{t_{k^+}}
  &= n_\text{eff}\bigl(\vr{\text{wye}}{t^\phi_k} - \vr{\text{wye}}{n_\text{wye}}\bigr)
     - \bigl(R^\text{w}_x c^{r,x}_{x,\text{wye},k} - X^\text{w}_x c^{i,x}_{x,\text{wye},k}\bigr)
     - n_\text{eff}\bigl(R^\text{d}_x c^{r,x}_{x,\text{del},k} - X^\text{d}_x c^{i,x}_{x,\text{del},k}\bigr),
  \label{eq:yd-v-r}\\
  \vi{\text{del}}{t_k} - \vi{\text{del}}{t_{k^+}}
  &= n_\text{eff}\bigl(\vi{\text{wye}}{t^\phi_k} - \vi{\text{wye}}{n_\text{wye}}\bigr)
     - \bigl(R^\text{w}_x c^{i,x}_{x,\text{wye},k} + X^\text{w}_x c^{r,x}_{x,\text{wye},k}\bigr)
     - n_\text{eff}\bigl(R^\text{d}_x c^{i,x}_{x,\text{del},k} + X^\text{d}_x c^{r,x}_{x,\text{del},k}\bigr),
  \label{eq:yd-v-i}
\end{align}
%
where $n_\text{wye}$ denotes the wye-side neutral terminal (if present; otherwise
the wye-side neutral voltage is taken to be zero).  When all impedance fields
are zero these collapse to the ideal transform.  The no-load shunt contributes
a current $Y_0 V$ at the from-side phase terminals, entering KCL exactly as the
standalone shunt of~\S\ref{sec:shunts}.

\paragraph{Current constraints.}
The current transform is the transpose of the voltage transform, ensuring power
conservation.  Let $k^- = ((k-2)\bmod n_\phi)+1$:
%
\begin{align}
  n_\text{eff}\,c^{r,x}_{x,\text{del},k}
  &= c^{r,x}_{x,\text{wye},k} - c^{r,x}_{x,\text{wye},k^-},
  \label{eq:yd-i-r}\\
  n_\text{eff}\,c^{i,x}_{x,\text{del},k}
  &= c^{i,x}_{x,\text{wye},k} - c^{i,x}_{x,\text{wye},k^-}.
  \label{eq:yd-i-i}
\end{align}

\paragraph{Neutral KCL at the transformer star point.}
When the wye winding has an explicit neutral terminal:
%
\begin{equation}
  c^{r,x}_{x,\text{wye},n} + \sum_{k=1}^{n_\phi} c^{r,x}_{x,\text{wye},k} = 0,
  \qquad
  c^{i,x}_{x,\text{wye},n} + \sum_{k=1}^{n_\phi} c^{i,x}_{x,\text{wye},k} = 0.
  \label{eq:yd-neutral}
\end{equation}

\subsection{Kirchhoff's Current Law}

KCL is enforced at every ungrounded terminal.  Each component function
accumulates its signed current contribution into per-terminal accumulators
$(\kappa^r_{b,t},\,\kappa^i_{b,t})$ using the sign convention
\emph{positive current flows into the bus}.  The final KCL constraints are:
%
\begin{equation}
  \kappa^r_{b,t} = 0,\qquad
  \kappa^i_{b,t} = 0,\qquad
  \forall\,(b,t)\notin\mc{G}_\text{nd}.
  \label{eq:kcl}
\end{equation}

Table~\ref{tab:kcl} summarises the contribution of each component type to the
KCL accumulator at a given terminal.

\begin{table}[h!]
\centering
\caption{KCL contributions (real part; imaginary part analogous).}
\label{tab:kcl}
\renewcommand{\arraystretch}{1.3}
\begin{tabular}{@{}llll@{}}
\toprule
Component & Terminal & Contribution & Sign rationale \\
\midrule
Voltage source $v$ & $t^\text{src}_k$ & $+c^{r,v}_{v,k}$ & Source injects into bus \\
Line $\ell$ & $t^\text{fr}_k$ (from-bus) & $-\bigl(c^r_{\ell,k} + I^{\text{sh},r}_k(b^\text{fr})\bigr)$ & Series + from-shunt leave \\
Line $\ell$ & $t^\text{to}_k$ (to-bus) & $-\bigl(\tilde{c}^r_{\ell,k} + I^{\text{sh},r}_k(b^\text{to})\bigr)$ & Series + to-shunt leave \\
Switch $s$ (closed) & $t^\text{fr}_k$ & $-c^{r,s}_{s,k}$ & Current leaves \\
Switch $s$ (closed) & $t^\text{to}_k$ & $+c^{r,s}_{s,k}$ & Current enters \\
Load $d$ (WYE) & phase $t^\phi_k$ & $-c^{r,d}_{d,k}$ & Load consumes current \\
Load $d$ (WYE) & neutral $n_b$ & $+c^{r,d}_{d,k}$ & Return path \\
Load $d$ (DELTA) & $t_k$ (positive) & $-c^{r,d}_{d,k}$ & Current leaves \\
Load $d$ (DELTA) & $t_{k^+}$ (negative) & $+c^{r,d}_{d,k}$ & Current returns \\
Generator $g$ (WYE) & phase $t^\phi_k$ & $+c^{r,g}_{g,k}$ & Generator injects \\
Generator $g$ (WYE) & neutral $n_b$ & $-c^{r,g}_{g,k}$ & Return path \\
IBR $n$ (FOUR\_LEG) & phase $t^\phi_k$ & $+c^{r,n}_{n,k}$ & IBR injects \\
IBR $n$ (FOUR\_LEG) & neutral $n_b$ & $-c^{r,n}_{n,k}$ & Return path \\
Transformer $x$ & any terminal $t$ & $-c^{r,x}_{x,\sigma,k}$ & All winding currents leave \\
\bottomrule
\end{tabular}
\end{table}

% ============================================================
\section{Feasibility Relaxation}
\label{sec:feasibility}

The standard OPF may be infeasible when the constant-power load and generation
specifications cannot be reconciled with the network's physical constraints.  To
enable diagnostics on such cases, a \emph{feasibility-relaxed} variant
$\texttt{solve\_feasibility\_opf}$ is provided.

\subsection{Elastic Slack Currents}

An elastic slack current $(c^{r,\varepsilon}_{b,t},\, c^{i,\varepsilon}_{b,t})$
is added at every ungrounded, non-source terminal.  These variables enter KCL
directly:
%
\begin{equation}
  \kappa^r_{b,t} + c^{r,\varepsilon}_{b,t} = 0,\qquad
  \kappa^i_{b,t} + c^{i,\varepsilon}_{b,t} = 0,
  \qquad \forall\,(b,t)\notin\mc{G}_\text{nd}\cup\mc{G}_\text{src}.
\end{equation}
%
Since the slacks are unconstrained, the modified KCL system is always feasible.

\subsection{Objective}

The cost objective is replaced by the $\ell_2^2$ norm of all slack injections:
%
\begin{equation}
  \min\;\sum_{(b,t)\notin\mc{G}_\text{nd}\cup\mc{G}_\text{src}}
    \Bigl[
      \bigl(c^{r,\varepsilon}_{b,t}\bigr)^2
      + \bigl(c^{i,\varepsilon}_{b,t}\bigr)^2
    \Bigr].
\end{equation}
%
All other constraints (KVL, power equations, transformer coupling) are
unchanged.  Voltage bounds are \emph{not} enforced as hard constraints in the
feasibility OPF; they are evaluated post-solve and reported separately.

\subsection{Interpretation}

A zero slack solution ($c^\varepsilon = 0$) certifies that the network is
physically feasible under its loading conditions.  Non-zero slacks at terminal
$(b,t)$ indicate that KCL cannot be satisfied there without external
intervention.  The magnitude and location of the slacks are used by
$\texttt{diagnose\_infeasibility}$ to rank and classify the root cause
(overloaded feeder, floating neutral, isolated bus, transformer mismatch, etc.).

% ============================================================
\section{Warm-Start Strategy}

The IVR formulation is an NLP with multiple local minima.  Without
initialisation, Ipopt typically converges to the degenerate $v = 0$ solution.
The following start values are applied before solving:

For the standard OPF:
\begin{equation}
  v^{r,0}_{b,t} = V^\text{src} \cos\theta_t,\quad
  v^{i,0}_{b,t} = V^\text{src} \sin\theta_t,
\end{equation}
where $V^\text{src}$ is the source voltage magnitude and
$\theta_t \in \{0,\,-2\pi/3,\,+2\pi/3,\,0\}$ for phase terminals
$\{$a, b, c, n$\}$.

For the feasibility OPF, a BFS voltage propagation through the transformer
tree assigns per-bus nominal voltages $V^\text{nom}_b$, and the start value
is $V^\text{nom}_b\cos\theta_t$ per bus.  This correctly initialises LV buses
at $\approx 230$\,V rather than at the source voltage ($\approx 6350$\,V for
an 11\,kV MV feeder), preventing convergence to the high-voltage local minimum
that arises without voltage bounds.

% ============================================================
\section{Summary of Constraint Counts}

For a network with $B$ buses, average $T$ terminals per bus, $L$ lines with
$n_\ell$ conductors each, $D$ wye-connected loads, $G$ generators, and one
voltage source, the model has approximately:
%
\begin{center}
\renewcommand{\arraystretch}{1.2}
\begin{tabular}{@{}lr@{}}
\toprule
Variable family & Count (real) \\
\midrule
Bus voltages $(v^r, v^i)$ & $2 B T$ \\
Line currents (fr + to) & $4 \sum_\ell n_\ell$ \\
Load currents & $2 \sum_d n^\phi_d$ \\
Generator currents & $2 \sum_g n^\phi_g$ \\
Source slack currents & $2 n_v$ \\
\midrule
Equality constraints & \\
\quad Grounded terminals & $2|\mc{G}_\text{nd}|$ \\
\quad Source voltage fixing & $2 n_v$ \\
\quad KVL (lines, real + imaginary) & $4\sum_\ell n_\ell$ \\
\quad Series current balance & $4\sum_\ell n_\ell$ \\
\quad Load P/Q equations & $4\sum_d n^\phi_d$ \\
\quad Transformer equations & linear in winding count \\
\quad KCL & $2(BT - |\mc{G}_\text{nd}|)$ \\
Inequality constraints & \\
\quad Voltage bounds & $\leq 2\sum_b |\mc{T}^\phi_b|$ \\
\quad Current limits & $\leq 2\sum_\ell n_\ell$ \\
\quad Generator P/Q bounds & $4\sum_g n^\phi_g$ \\
\bottomrule
\end{tabular}
\end{center}

% ============================================================
\section{Relationship to PowerModelsDistribution}

The formulation follows the \texttt{IVRENPowerModel} in
PowerModelsDistribution~\cite{PowerModelsDistribution2022}, with these
differences:
%
\begin{itemize}
  \item All variables and parameters are in \textbf{SI units} (V, A, W, var,
        $\Omega$) rather than per-unit.

  \item The neutral terminal is identified by the \texttt{neutral\_terminal}
        bus field or the ``n''/``N'' naming convention, rather than by a fixed
        conductor index.

  \item Network data is read directly from the BMOPF JSON schema rather than
        the PMD MATHEMATICAL model dict.

  \item The feasibility relaxation and infeasibility diagnostics are
        BMOPFTools-specific extensions not present in PMD.
\end{itemize}

% ============================================================
\bibliographystyle{IEEEtran}
\begin{thebibliography}{9}

\bibitem{PowerModelsDistribution2022}
D.~Fobes, S.~ Claeys, F.~Geth, and C.~Coffrin,
``PowerModelsDistribution.jl: An Open-Source Framework for Exploring
  Distribution Power Flow Formulations,''
\emph{Electric Power Systems Research}, vol.~189, p.~106664, 2020.

\end{thebibliography}

\end{document}
